\documentclass[11pt]{article}
\usepackage[utf8]{inputenc}
\usepackage[T1]{fontenc}
\usepackage[polish]{babel}
\usepackage{amsmath,amssymb}
\usepackage{booktabs}
\usepackage{geometry}
\geometry{margin=2.5cm}
\usepackage{hyperref}
\usepackage{pgfplots}
\pgfplotsset{compat=1.17}

\title{Kompozycja małych ekspertów muzycznych:\\ baseline'y, bezstratny stitch i uczciwy wynik negatywny}
\author{Arkadiusz Słota \\ kolektyw \textbf{Slayer} \\ \texttt{github.com/slayerlabs/micro-models}}
\date{czerwiec 2026}

\begin{document}
\maketitle

\begin{abstract}
Trenujemy cztery modele GPT ($\sim$0{,}8 mln parametrów) \emph{od zera} na monofonicznej muzyce w notacji ABC
(jig irlandzki, sopran chorałów Bacha, walc, reel). Mimo małego rozmiaru uczą się struktury muzycznej --- metrum,
sygnatur tonacji, kadencji --- z samej predykcji następnego znaku, bez wpisanej w kod teorii muzyki. Następnie
badamy, czy takich małych ekspertów da się \emph{składać}. Trzy wnioski: (i) mechanizm stitchu jest
\textbf{bezstratny} --- wytrenowany liniowy mapper na styku pośrednim odtwarza baseline; (ii) płaski
\textbf{ensemble} dwóch ekspertów bije każdy pojedynczy model na zbalansowanym zadaniu mieszanym; ale
(iii) \textbf{stitch reprezentacji} (liniowy) tego ensemble \emph{nie bije}. Raportujemy ten wynik negatywny
otwarcie i wskazujemy warunki --- wymuszony wspólny kontrakt, bogatsze mappery --- w których kompozycja na
poziomie reprezentacji mogłaby wygrać.
\end{abstract}

\section{Motywacja}
Interesują nas \emph{mikro-modele}: sieci dość małe, by zrozumieć je od początku do końca i wytrenować w minuty.
Pytanie przewodnie: czy wielu małych, poprawnych ekspertów da się \textbf{złożyć} w coś większego --- taniej i
w sposób modularny --- zamiast trenować jeden duży model. Predykcja następnego tokena to wspólny cel od n-gramu
po transformer; my pracujemy na jego najmniejszym praktycznym końcu. Zgodnie z zasadami kolektywu Slayer
(twardy pomiar, jawny koszt, brak benchmaxxingu) raportujemy też wyniki negatywne.

\section{Modele i dane}
Wszyscy eksperci dzielą jedną architekturę: transformer dekoderowy, \textbf{4 warstwy, 4 głowice,
$d_\text{model}=128$, okno 128 znaków}, poziom znaku (char-level). Trening: strata cross-entropy, AdamW
($\text{lr}=3\cdot10^{-4}$), warmup + cosine, $\sim$2000 kroków, CPU (minuty). Dane: melodie folkowe z
thesession.org (ABC) oraz soprany chorałów Bacha przez \texttt{music21}; surowych danych nie redystrybuujemy
(odtwarzalne ze skryptów).

\begin{table}[h]
\centering
\begin{tabular}{llrr}
\toprule
Ekspert & Styl / render & Dane (znaki) & Val ppl \\
\midrule
jig    & jig irlandzki 6/8           & 2{,}45 mln & 3{,}80 \\
Bach   & sopran chorału (barok)      & 79 tys.    & 2{,}09$^\ast$ \\
walc   & liryczny 3/4 $\to$ fortepian & 789 tys.  & $\sim$4{,}4 \\
reel   & żwawy fiddle 4/4 $\to$ skrzypce & 3{,}90 mln & $\sim$4{,}9 \\
\bottomrule
\end{tabular}
\caption{Eksperci. $^\ast$ppl Bacha nieporównywalne 1:1 (mniejszy słownik + bardzo powtarzalne dane).}
\end{table}

Modele uczą się realnej struktury: generują poprawne ABC, respektują metrum i sygnatury tonacji (np. $^\sharp$F
w G-dur pojawia się samo, bez błędnych krzyżyków) --- wyłącznie z predykcji następnego znaku.

\section{Eksperymenty kompozycji}
\paragraph{E0 --- self-stitch (mechanizm).} Wstawiamy liniowy mapper $g\in\mathbb{R}^{128\times128}$ na styku po
bloku 2 tego samego modelu i zamrażamy resztę. $g=$ identyczność daje \emph{dokładnie} baseline (poprawność
forwardu); $g$ losowy psuje model (ppl $122$), a trening samych 16 tys.\ parametrów $g$ odbudowuje go do
baseline (Δppl $+0{,}03$). \textbf{Waliduje to mechanizm, nie tezę kompozycji.}

\paragraph{Ensemble (płaskie ważenie).} Dwóch ekspertów o wspólnym słowniku; w każdym kroku
$\text{logits}=\alpha\,A + (1-\alpha)\,B$. Baseline do pobicia.

\paragraph{E1 --- stitch reprezentacji.} Front eksperta A $\times\, g \,\times$ back eksperta B (wspólny słownik),
trening tylko $g$. \emph{Lekcja metodologiczna:} pierwszy pomiar dał ,,stitch bije ensemble'' --- okazało się
artefaktem (held-out ze sklejonych korpusów wyszedł niemal czysty reel). Po zbalansowaniu (val = 50\% walc +
50\% reel) wynik się odwrócił. Setup ewaluacji jest częścią dowodu.

\section{Wyniki}
\begin{figure}[h]
\centering
\begin{tikzpicture}
\begin{axis}[
  ybar, width=11cm, height=6cm,
  symbolic x coords={walc, reel, ensemble, stitch},
  xtick=data, ylabel={val perplexity (niżej = lepiej)},
  nodes near coords, nodes near coords precision=2,
  ymin=0, ymax=7.2, bar width=22pt, enlarge x limits=0.22]
\addplot coordinates {(walc,6.20) (reel,5.87) (ensemble,5.15) (stitch,5.18)};
\end{axis}
\end{tikzpicture}
\caption{Zbalansowany held-out (walc + reel po równo). \textbf{Ensemble bije oba pojedyncze modele;
stitch reprezentacji $\approx$ ensemble (minimalnie gorszy).}}
\end{figure}

\noindent Dwa wnioski: \textbf{(1) ensemble (5{,}15) bije oba pojedyncze modele} (walc 6{,}20, reel 5{,}87) ---
łączenie ekspertów realnie pomaga na zadaniu mieszanym. \textbf{(2) Stitch (5{,}18) $\approx$ ensemble (5{,}15)}
--- liniowy stitch reprezentacji nie pobił prostego baseline'u.

\section{Dyskusja}
Wynik (2) jest negatywny, ale spodziewany: zrobiliśmy najtrudniejszy przypadek --- \emph{post-hoc}, mapper
\emph{liniowy}, \emph{bez} wymuszonego wspólnego kontraktu, z pominięciem testu wyrównania geometrii. Dwa modele
trenowane niezależnie mają różne geometrie reprezentacji; pojedynczy liniowy mapper na jednym styku nie wyrównuje
ich lepiej niż uśrednianie wyjść. Drogi do ewentualnej przewagi: (a) \textbf{wymuszony kontrakt} --- trenować obu
ekspertów przeciw wspólnej, zamrożonej głowicy od zera \cite{rebasin,zipit}; (b) \textbf{bogatszy mapper} ---
słownik cech / sparse autoencoder zamiast liniowego \cite{vqvae,sae}; (c) najpierw sprawdzić, czy niezależne
modele tej samej architektury są w ogóle liniowo zszywalne \cite{stitch15,bansal21}. Prosty baseline (ensemble
\cite{ensembles}) jest mocny --- złożoność musi dopiero zasłużyć na siebie.

\paragraph{Asymetria porównania.} Nasz stitch jest \emph{kierunkowy} (front A $\times\,g\,\times$ głowa B), a ensemble \emph{symetryczny} (miesza obie głowy). Uczciwszy test pytania „czy kompozycja reprezentacji bije mieszanie wyjścia'' to \emph{symetryczna} fuzja reprezentacji do wspólnej głowy --- nie unieważnia remisu, ale go niuansuje.

\paragraph{Pomiar konwergencji reprezentacji (CKA na mikro-skali).} Sprawdziliśmy wprost, czy niezależnie trenowane mikro-modele dzielą geometrię --- metrykami CKA \cite{cka} i mutual-kNN \cite{platonic}, jako \emph{funkcję skali} (Rys.~\ref{fig:cka}). \textbf{Tak: konwergencja jest wysoka już na $0{,}05$M} (CKA $\approx0{,}95$, daleko powyżej poziomu null $0{,}35$ dla modelu-vs-losowy). To \emph{obala} naszą pierwotną interpretację „różnych geometrii'' --- geometrie są wspólne, więc remis stitch$\approx$ensemble wynika z \textbf{redundancji} (mało komplementarnej informacji), nie z niezgodności; wspólna geometria otwiera za to darmową interoperacyjność przez reprezentacje względne \cite{relreps}. \emph{Niuans:} CKA saturuje ($\approx0{,}95$--$0{,}97$), więc trend ze skalą widać dopiero w czulszej metryce mutual-kNN ($0{,}71\to0{,}84$, monotonicznie i ponad szum seedów) --- zgodnie z kierunkiem hipotezy platońskiej \cite{platonic}; \emph{wybór metryki ma znaczenie}. \emph{Uwaga do dalszego routingu:} domeny rozłączne po metrum (6/8 vs 4/4) pozwalają routerowi „oszukać'' po jednym tokenie nagłówka --- testy routingu wymagają domen w \emph{tym samym} metrum.

\begin{figure}[h]
\centering
\begin{tikzpicture}
\begin{axis}[
  width=12cm, height=6.5cm, xmode=log, log basis x=10,
  xlabel={rozmiar modelu (mln parametrów)}, ylabel={podobieństwo reprezentacji},
  xtick={0.05,0.2,0.8,1.8}, xticklabels={0{,}05M,0{,}2M,0{,}8M,1{,}8M},
  ymin=0.3, ymax=1.02, legend pos=south east, grid=major, mark size=2.5pt]
\addplot+[mark=*, error bars/.cd, y dir=both, y explicit] coordinates {
  (0.05,0.952)+-(0,0.007) (0.2,0.964)+-(0,0.006) (0.8,0.973)+-(0,0.000) (1.8,0.971)+-(0,0.000)};
\addlegendentry{CKA}
\addplot+[mark=square*, error bars/.cd, y dir=both, y explicit] coordinates {
  (0.05,0.713)+-(0,0.018) (0.2,0.793)+-(0,0.011) (0.8,0.829)+-(0,0.005) (1.8,0.842)+-(0,0.002)};
\addlegendentry{mutual-kNN ($k{=}10$)}
\addplot[black, dashed, mark=none] coordinates {(0.05,0.35) (1.8,0.35)};
\addlegendentry{null (vs losowy)}
\end{axis}
\end{tikzpicture}
\caption{\textbf{Konwergencja reprezentacji niezależnie trenowanych mikro-modeli w funkcji rozmiaru.} Podobieństwo \emph{same-task} (te same dane --- jigi ABC; różny seed) między wyjściami bloków, uśrednione po warstwach, mierzone CKA \cite{cka} i mutual-kNN ($k{=}10$; metryka z \cite{platonic}); słupki błędu $=\pm$std po 3 parach seedów na każdą skalę. Linia przerywana: poziom null (model wytrenowany vs losowy, CKA$\approx0{,}35$). \emph{Wniosek:} konwergencja jest wysoka już na $0{,}05$M; mutual-kNN rośnie monotonicznie i ponad szum seedów wraz ze skalą ($0{,}71\to0{,}84$, malejące przyrosty), zgodnie z kierunkiem hipotezy platońskiej, podczas gdy CKA saturuje ($\approx0{,}95$--$0{,}97$) i jest zbyt zgrubny, by ten trend rozdzielić --- wybór metryki ma znaczenie na mikro-skali.}
\label{fig:cka}
\end{figure}

\section{Wnioski}
Pokazaliśmy, że (i) malutki char-LM uczy się struktury muzycznej z samej predykcji następnego znaku oraz (ii)
\emph{ensemble} małych ekspertów bije pojedynczy model na zadaniu mieszanym --- kompozycja ma wartość. Nie
pokazaliśmy (jeszcze), że \emph{stitch na poziomie reprezentacji} bije ten ensemble; ten wynik negatywny wskazuje
konkretne, sfalsyfikowalne kierunki dalej. Lekcję metodologiczną traktujemy jako \emph{główną}, nie poboczną:
\textbf{setup ewaluacji jest częścią dowodu} --- pierwszy pomiar E1 dał pozorne „stitch bije ensemble'', okazał się
artefaktem skażonego held-outu (sklejone korpusy $\to$ walidacja niemal czysty reel), i sami go obaliliśmy przed
publikacją. Cały kod, wagi i notatki badawcze (dialektyka teza$\leftrightarrow$antyteza$\to$synteza) są otwarte w repozytorium.

\begin{thebibliography}{9}
\bibitem{stitch15} K.~Lenc, A.~Vedaldi. \emph{Understanding image representations by measuring their equivariance and equivalence}. CVPR 2015.
\bibitem{bansal21} Y.~Bansal, P.~Nakkiran, B.~Barak. \emph{Revisiting Model Stitching to Compare Neural Representations}. NeurIPS 2021.
\bibitem{zipit} G.~Stoica i in. \emph{ZipIt! Merging Models from Different Tasks without Training}. ICLR 2024.
\bibitem{rebasin} S.~Ainsworth, J.~Hayase, S.~Srinivasa. \emph{Git Re-Basin}. ICLR 2023.
\bibitem{ensembles} B.~Lakshminarayanan, A.~Pritzel, C.~Blundell. \emph{Simple and Scalable Predictive Uncertainty Estimation using Deep Ensembles}. NeurIPS 2017.
\bibitem{vqvae} A.~van den Oord, O.~Vinyals, K.~Kavukcuoglu. \emph{Neural Discrete Representation Learning} (VQ-VAE). NeurIPS 2017.
\bibitem{sae} T.~Bricken i in. \emph{Towards Monosemanticity} (sparse autoencoders), 2023; H.~Cunningham i in., arXiv:2309.08600.
\bibitem{nplm} Y.~Bengio i in. \emph{A Neural Probabilistic Language Model}. JMLR 2003.
\bibitem{knnlm} U.~Khandelwal i in. \emph{Generalization through Memorization: Nearest Neighbor Language Models}. ICLR 2020.
\bibitem{platonic} M.~Huh, B.~Cheung, T.~Wang, P.~Isola. \emph{Position: The Platonic Representation Hypothesis}. ICML 2024. arXiv:2405.07987.
\bibitem{relreps} L.~Moschella i in. \emph{Relative representations enable zero-shot latent space communication}. ICLR 2023. arXiv:2209.15430.
\bibitem{cka} S.~Kornblith, M.~Norouzi, H.~Lee, G.~Hinton. \emph{Similarity of Neural Network Representations Revisited}. ICML 2019. arXiv:1905.00414.
\end{thebibliography}

\end{document}
